% TOP_PROBLEM: {"id":30007656,"problem_number":"LOCAL-30007656","title":"Targeting need versus potential impact","category_id":21,"category":{"id":21,"name":"economics","display_name":"Economics","description":"Open economic theory statements, published research agendas, and proposed scoped specifications; question provenance and historical priority are qualified per record.","slug":"economics","order_index":21,"created_at":"2026-10-08T00:00:00Z"},"status":"research_question","external_url":null,"legacy_ids":[],"published":false,"statement":"In the explicitly specified setting: How should antipoverty programs balance reaching the most deprived against reaching those who gain most, while limiting exclusion, manipulation, administrative costs, and political capture? The mathematical statement above fixes the target, assumptions and scope of this catalogue formulation.","economics":{"original_id":145,"field":"Development and poverty","kind":"design","formalization_basis":"adapted_research_question","source_status":"research_agenda","priority_status":"earliest_found","priority_note":"Earliest directly inspected qualifying passage in this audit; earlier origins were not established. A review or research-agenda statement does not imply original priority.","earliest_verified_reference":{"authors":["Johannes Haushofer","Paul Niehaus","Carlos Paramo","Edward Miguel","Michael W. Walker"],"title":"Targeting Impact versus Deprivation","year":2025,"url":"https://emiguel.econ.berkeley.edu/wordpress/wp-content/uploads/2026/01/haushofer-et-al-2025-targeting-impact-versus-deprivation.pdf","doi":"10.1257/aer.20221650","locator":"Section VI, Conclusion, printed pp. 1971–1972, first two caveats and final paragraph","evidence_summary":"The authors explicitly identify cross-context and long-horizon targeting uncertainties and call for extending the approach; strategic reporting is our additional specification."},"current_reference":{"authors":["Johannes Haushofer","Paul Niehaus","Carlos Paramo","Edward Miguel","Michael W. Walker"],"title":"Targeting Impact versus Deprivation","year":2025,"url":"https://emiguel.econ.berkeley.edu/wordpress/wp-content/uploads/2026/01/haushofer-et-al-2025-targeting-impact-versus-deprivation.pdf","doi":"10.1257/aer.20221650","locator":"Section VI, Conclusion, printed pp. 1971–1972, first two caveats and final paragraph","evidence_summary":"The authors explicitly identify cross-context and long-horizon targeting uncertainties and call for extending the approach; strategic reporting is our additional specification."},"formalization_note":"The cited passage establishes a research direction, not this exact mathematical statement. The notation, population, policy menu, assumptions, and estimands are an original adaptation for this catalogue; no universal unsolved theorem or source priority is claimed."},"statusQualification":"Published question or research agenda verified at the cited locator. The mathematical specification is adapted for this catalogue and includes additional assumptions; source publication does not establish first-ever priority or certify this exact model remains unsolved. The cited passage establishes a research direction, not this exact mathematical statement. The notation, population, policy menu, assumptions, and estimands are an original adaptation for this catalogue; no universal unsolved theorem or source priority is claimed.","statusReviewedAt":"2026-10-08"}
% FIRST_POSTED: 2026-10-08
% ENTRY_KIND: statement_only
% !TeX program = lualatex
\documentclass[11pt]{article}
\usepackage{fontspec}
\usepackage[a4paper,margin=25mm]{geometry}
\usepackage{amsmath,amssymb,mathtools}
\usepackage{xurl}
\usepackage[hidelinks]{hyperref}
\newcommand{\sref}[2]{\hyperlink{source:#1:#2}{[S#2]}}
\setlength{\emergencystretch}{3em}
\begin{document}
% problemId:30007656
\section*{Targeting need versus potential impact}\label{30007656}
\subsection{Definitions and mathematical statement}
Fix a target population, an assistance program, and a fixed per-household fiscal budget \(B\). Let \(X_i\) be baseline information legally and practically available for household \(i\), \(c_i\) the full delivery and targeting cost, and \(Y_i(d)\) a declared welfare outcome under eligibility \(d\in\{0,1\}\). Define need \(n_i\) with a predeclared deprivation measure and expected impact \(\tau(x)=E[Y_i(1)-Y_i(0)\mid X_i=x]\). A targeting rule \(\pi(x)\in[0,1]\) assigns eligibility probabilities. Define weighted gain \(g(x)=E[\omega(n_i)\{Y_i(1)-Y_i(0)\}\mid X_i=x]\). Characterize rules maximizing \(E[g(X_i)\pi(X_i)]\) subject to \(E[c_i\pi(X_i)]\le B\) and explicit minimum-coverage or equity constraints. The weight \(\omega\) is a declared social choice, not an estimated fact. Determine how the frontier changes when households manipulate reports, communities select recipients, or implementation creates exclusion and political capture. Randomized program assignment and independent measurement are needed to estimate impact; a deprivation predictor is not automatically an impact predictor. Validate rules in new settings and at scale, including administrative costs. Existing targeting studies answer important specific comparisons. Their agenda calls for other contexts, interventions, and longer horizons. The objective and strategic-reporting extensions here adapt that agenda; they are not a verbatim published formal problem.

\par\textbf{Formulation and scope.} The cited passage establishes a research direction, not this exact mathematical statement. The notation, population, policy menu, assumptions, and estimands are an original adaptation for this catalogue; no universal unsolved theorem or source priority is claimed.

\par\textbf{Open-question provenance.} An explicit question or research agenda was verified at the source locator. This does not by itself certify that every later version remains unresolved \sref{30007656}{1}.

\par\textbf{Historical priority.} Earliest directly inspected qualifying passage in this audit; earlier origins were not established. A review or research-agenda statement does not imply original priority.

\subsection{Short English statement}
In the explicitly specified setting: How should antipoverty programs balance reaching the most deprived against reaching those who gain most, while limiting exclusion, manipulation, administrative costs, and political capture? The mathematical statement above fixes the target, assumptions and scope of this catalogue formulation.

\subsection{Sources}
\begin{itemize}\raggedright
\item[\textnormal{[S1]}]\hypertarget{source:30007656:1}{} Johannes Haushofer, Paul Niehaus, Carlos Paramo, Edward Miguel, Michael W. Walker. 2025. Targeting Impact versus Deprivation. Section VI, Conclusion, printed pp. 1971–1972, first two caveats and final paragraph.
\url{https://emiguel.econ.berkeley.edu/wordpress/wp-content/uploads/2026/01/haushofer-et-al-2025-targeting-impact-versus-deprivation.pdf}.
DOI: 10.1257/aer.20221650.
{\small\textit{Earliest verified question/agenda reference; historical priority qualified below. The authors explicitly identify cross-context and long-horizon targeting uncertainties and call for extending the approach; strategic reporting is our additional specification.}}
\end{itemize}
\end{document}
